\documentclass[10pt]{article}
\usepackage[margin=1.5cm]{geometry}
\usepackage{amsmath}
\usepackage[spanish]{babel}
\usepackage{parskip}
\usepackage[most]{tcolorbox}
\pagestyle{empty}

\tcbset{
  mybox/.style={
    enhanced,
    skin=enhancedmiddle,
    colback=white,
    colframe=black,
    boxrule=0.4pt,
    sharp corners,
    width=0.48\textwidth,
    height=0.23\textheight,
    boxsep=2pt,
    left=4pt,
    right=4pt,
    top=4pt,
    bottom=4pt,
    borderline={0.4pt}{0pt}{dash pattern=on 3pt off 2pt},
  }
}


\begin{document}

% Evaluación 1
\begin{tcolorbox}[mybox]
\textbf{Escuela: 4-117 Ejército de los Andes} \\
\textbf{Curso: 5º 3ra División} \hfill Nombre: \hrulefill

\vspace{1mm}
\textbf{Evaluación 1 – Mecánica de los Fluidos} \\
\textbf{Duración: 80 minutos}

\begin{enumerate}
\item Calcular la presión absoluta a $7{,}6\,\text{m}$ de profundidad en agua.
\item ¿Qué profundidad de agua produce $2{,}066\,\text{kg/cm}^2$ de presión?
\item Se aplica $F=40\,\text{N}$ sobre un émbolo de $A=4\,\text{cm}^2$. ¿Cuál es la presión generada?
\item En una prensa con $A_1=0{,}001\,\text{m}^2$, $A_2=0{,}1\,\text{m}^2$ y $F_1=100\,\text{N}$, calcular $F_2$.
\item ¿Qué fuerza debe aplicarse en un émbolo de $r=0{,}05\,\text{m}$ para levantar $P=1000\,\text{N}$ con un émbolo de $r=0{,}25\,\text{m}$?
\end{enumerate}
\end{tcolorbox}
\hfill
% Evaluación 2
\begin{tcolorbox}[mybox]
\textbf{Escuela: 4-117 Ejército de los Andes} \\
\textbf{Curso: 5º 3ra División} \hfill Nombre: \hrulefill

\vspace{1mm}
\textbf{Evaluación 2 – Mecánica de los Fluidos} \\
\textbf{Duración: 80 minutos}


\begin{enumerate}
\item Determinar la presión a $10\,\text{m}$ de profundidad en agua.
\item ¿Qué presión transmite una fuerza de $60\,\text{N}$ sobre un área de $6\,\text{cm}^2$?
\item En una prensa, el émbolo mayor tiene radio $r=0{,}3\,\text{m}$ y el menor $r=0{,}05\,\text{m}$. Se aplica $F_1=120\,\text{N}$. Calcular $F_2$.
\item ¿Qué masa puede levantarse si en el émbolo menor ($r=4\,\text{cm}$) se colocan $6\,\text{kg}$ y el émbolo mayor tiene radio triple?
\item ¿Qué altura debe tener una columna de agua para generar $1{,}5\times10^5\,\text{Pa}$?
\end{enumerate}
.
\end{tcolorbox}

\end{document}
